% Solution: Tadelis (2013) Exercise 15.1 "Equilibrium Selection". Input by \solitem.
% Rewritten and deepened from ../pratica-lecture-08-solutions/sol-tadelis-15-1.tex. Check: check_G_L07.py.
\begin{exsol}{Tadelis Exercise 15.1}{314}{Foundations}{1}
\asked (a) All Bayesian Nash equilibria (BNE) of the tree in Figure~15.5; (b) which of them are
perfect Bayesian equilibria (PBE), and why.

\method A BNE of an extensive-form game is a Nash equilibrium of its normal form, so write the
$3\times3$ matrix and find its equilibria. For the PBE, give player~2 a belief $\mu = \Pr(x_2)$ at his
information set $\{x_2, x_3\}$ and require his action to be a best reply to some belief: Bayes' rule
if the set is reached, any belief if it is not (Tadelis \S15.1--15.2, pp.~304--312). Slides:
Lecture~7, frames~30--31 (the BNE that relies on a non-credible move), 33 (system of beliefs),
34--35 (Requirements~1--4) and 36 (definition of PBE).

\solution
\textbf{Reading the tree.} Player~1 chooses $A$, $B$ or $C$; $C$ ends the game with $(1, 3)$. After $A$
(node $x_2$) or $B$ (node $x_3$) player~2 chooses $L$, $M$ or $R$ without knowing which of the two nodes
he is at. Payoffs $(u_1, u_2)$: after $A$: $L\,(0,4)$, $M\,(0,1)$, $R\,(4,0)$; after $B$: $L\,(4,0)$,
$M\,(0,1)$, $R\,(0,4)$.

\textbf{Normal form} (best replies underlined):
\begin{center}
\begin{tabular}{r|c|c|c|}
\multicolumn{1}{r}{} & \multicolumn{1}{c}{$L$} & \multicolumn{1}{c}{$M$} & \multicolumn{1}{c}{$R$}\\ \cline{2-4}
$A$ & $0,\ \br{4}$ & $0,\ 1$ & $\br{4},\ 0$ \\ \cline{2-4}
$B$ & $\br{4},\ 0$ & $0,\ 1$ & $0,\ \br{4}$ \\ \cline{2-4}
$C$ & $1,\ \br{3}$ & $\br{1},\ \br{3}$ & $1,\ \br{3}$ \\ \cline{2-4}
\end{tabular}
\end{center}
Against $L$ player~1's best reply is $B$ (4); against $M$ it is $C$ (1 beats 0); against $R$ it is $A$.
Against $A$ player~2's best reply is $L$, against $B$ it is $R$, and against $C$ he is indifferent.

\textbf{(a) Pure BNE.} The only cell with both payoffs underlined is $(C, M)$.

\emph{Mixed BNE.} Let player~1 play $(a, b, c)$ and player~2 $(l, m, r)$. If $a + b > 0$ the information
set is reached, with belief $\mu = a/(a+b)$, and
\[
  u_2(L) = 4\mu, \qquad u_2(M) = 1, \qquad u_2(R) = 4(1 - \mu).
\]
Since $\max\{4\mu, 4(1-\mu)\} \ge 2 > 1$, $M$ is never a best reply when the set is reached: $m = 0$.
\begin{itemize}
  \item \emph{$c = 1$.} Player~1 gets $1$; deviating to $A$ gives $4r$ and to $B$ gives $4l$. Equilibrium
  iff $l \le \frac14$ and $r \le \frac14$ (so $m \ge \frac12$); player~2's choice is unconstrained because
  his set is not reached. This family contains $(C, M)$.
  \item \emph{$a, b > 0$.} Then $m = 0$; player~1 mixes $A$ and $B$ only if $4r = 4l$, so $l = r = \frac12$,
  and player~2 mixes $L$ and $R$ only if $4\mu = 4(1-\mu)$, so $\mu = \frac12$, i.e.\ $a = b$. Each of $A$,
  $B$ pays $2 > 1$, so $c = 0$: the profile $\bigl(\frac12 A + \frac12 B,\ \frac12 L + \frac12 R\bigr)$.
  \item \emph{Only one of $a, b$ positive} (say $a > 0 = b$): $\mu = 1$, player~2 plays $L$, and then $B$
  pays $4 > 0$, so player~1 would not play $A$. Mixing $C$ with $A$ and/or $B$ needs $4r = 1$ or
  $4l = 1$ with $m = 0$ and the set reached, which the first two cases rule out.
\end{itemize}

\textbf{(b) PBE.} \emph{The family with $c = 1$.} The information set is off the path, so Bayes' rule says
nothing and $\mu$ may be anything (Requirement~3). But Requirement~4 still binds: player~2 must best
respond to \emph{some} belief. $M$ is a best reply to none ($\max\{4\mu, 4(1-\mu)\} \ge 2 > 1$), so
$m = 0$, and then $l + r = 1$ contradicts $l, r \le \frac14$. No member of the family is a PBE; in
particular $(C, M)$ is a BNE that is not a PBE. \emph{The mixed profile.} The set is reached with
$\mu = \frac12$ by Bayes' rule; then $L$ and $R$ both give $2 > 1$, so mixing them is sequentially
rational, and player~1's $A$ and $B$ each give $2 > 1$ from $C$. It is a PBE with $\mu = \frac12$.

\interp $(C, M)$ is the Lecture~7 entry-game problem (frame~30) in miniature: player~1 is deterred from
entering the information set by a reply, $M$, that player~2 would never choose once there. Subgame
perfection cannot remove it, because the only proper subgame is the whole game; asking for a belief at
the information set and a best reply to it does.

\verify \texttt{check\_G\_L07.py} finds the pure equilibrium of the matrix, enumerates every mixed
profile with probabilities in multiples of $\frac1{12}$ (the equilibria found are exactly the two families
above), tests the boundary $l = \frac14$ and $l = \frac14 + \frac1{1000}$, and checks that $M$ is never a best
reply on a belief grid. Output: ``Tadelis 15.1: pure BNE (C,M); BNE = \{C with l,r <= 1/4\} +
(A/2+B/2, L/2+R/2); only the mixed one is a PBE (mu = 1/2)''. Tadelis prints no answer.

\mistakes
\begin{itemize}
  \item Calling $(C, M)$ a PBE ``because beliefs off the path are free''. Free beliefs do not license any
  action: the action must be optimal for at least one belief, and $M$ is optimal for none.
  \item Looking only for pure equilibria and concluding the game has no PBE. The only PBE is mixed.
  \item Checking subgame perfection instead: there is no proper subgame below the root (the information
  set has two nodes), so every BNE is an SPE here and SPE removes nothing.
  \item Forgetting that the BNE with $c = 1$ form a whole family ($l, r \le \frac14$), not just $(C, M)$.
\end{itemize}

\answer{(a) BNE: player~1 plays $C$ and player~2 plays $L$ and $R$ each with probability at most
$\frac14$ (pure case $(C, M)$); and $\bigl(\frac12 A + \frac12 B,\ \frac12 L + \frac12 R\bigr)$. (b) Only the
mixed profile, with belief $\mu = \frac12$, is a PBE. $(C, M)$ is not, because $M$ is never a best reply
at player~2's information set, whatever he believes.}
\end{exsol}
